\section{A Small Certificate Kernel}
\label{sec:framework}

This section describes the implemented interfaces in
\code{prototype/interface/GuardInterface.v}. State, code, checks, and
observations are parameters. The kernel neither interprets C expressions nor
extracts assumptions from arbitrary source/candidate pairs.

\subsection{Host Semantics and Private Check State}

A \code{guard_host} supplies relations $\run{c}{s}{o}$ for code execution and
$\checkrun{g}{s}{(b,s')}$ for check execution, together with a safety predicate
and a guarded-choice constructor $\operatorname{select}(g,t,f)$. Its dispatch law
is exact:
\begin{equation}
\run{\operatorname{select}(g,t,f)}{s}{o}
\ \Longleftrightarrow\
\exists b,s'.\ \checkrun{g}{s}{(b,s')}\ \land\
\run{(\mathbf{if}\ b\ \mathbf{then}\ t\ \mathbf{else}\ f)}{s'}{o}.
\label{eq:dispatch}
\end{equation}
The language instance proves this law for its declared observations. It must
account for observable check effects and any exceptional or divergent behavior
it claims to include. The constructor need not be a single syntactic conditional.

Checks may alter private temporaries. Consequently, the interface retains both
the original entry $s$ and the checked state $s'$. Accepted and refused entry
relations $R_+(s,s')$ and $R_-(s,s')$ state what the selected branch may rely on.
For a read-only frontend, these relations can be equality. For a private cursor
scan, they instead express preservation of public temporaries and memory.

\subsection{Guard and Transformation Certificates}

A \code{guard_certificate} is parameterized by an entry domain $D$, premise $P$,
the two entry relations, and the actual check $g$. It requires:
\begin{align}
D(s)&\Rightarrow\operatorname{safe}(g,s),\label{eq:safety}\\
D(s)&\Rightarrow\exists b,s'.\ \checkrun{g}{s}{(b,s')},\label{eq:available}\\
D(s)\land\checkrun{g}{s}{(b,s')}&\Rightarrow
\begin{cases}P(s)\land R_+(s,s')&b=\mathit{true},\\R_-(s,s')&b=\mathit{false}.
\end{cases}\label{eq:sound}
\end{align}
The premise is anchored at the original entry. Refusal establishes a valid
fallback entry, not $\neg P$. A conservative condition is therefore permitted.
Availability is existential; it is not a general check termination theorem.
Safety is a host predicate whose intended interpretation the language must
justify separately.

A \code{conditional_certificate} proves target-to-source refinement for both
branches. Let $o_t\preceq o_s$ be the supplied observation relation. Its accepted
clause is
\begin{equation}
D(s)\land P(s)\land R_+(s,s')\land\run{T}{s'}{o_t}
\ \Rightarrow\ \exists o_s.\ \run{S}{s}{o_s}\land o_t\preceq o_s.
\label{eq:conditional}
\end{equation}
The refused clause has $R_-$ and the fallback $F$, without requiring $P$.
When $F=S$, private-state transport still needs proof; syntactic reuse of the
source does not imply that it behaves correctly from $s'$.

\begin{theorem}[Local guarded refinement]
Let the guard and conditional certificates share a domain, premise, entry
relations, and branches. Every behavior of
$\operatorname{select}(g,T,F)$ from an entry satisfying $D$ is related to a
behavior of $S$ from that entry.
\end{theorem}
\begin{proof}
Apply the dispatch law to obtain the actual check result and checked state.
Guard soundness yields the appropriate entry relation and, on acceptance, $P$.
Apply the corresponding branch refinement clause.
\end{proof}
This is \code{guardify_refinement}. Source-to-target behavior preservation is
provided independently by \code{preservation_certificate} and
\code{guardify_preservation}, which also use check availability. Neither theorem
silently supplies the other direction or upgrades a finite big-step host to a
divergence theorem.

\subsection{Condition Libraries and Assumption Derivation}

The logical certificate chain distinguishes four links:
\begin{center}
\begin{tabular}{ll}
$C_{\mathrm{opt}}$&$A$ suffices for candidate correctness;\\
$C_{\mathrm{derive}}$&entry property $B$ implies the local/model obligations $A$;\\
$C_{\mathrm{guard}}$&the actual check safely establishes $B$ on acceptance;\\
$C_{\mathrm{host}}$&concrete guarded choice implements Eq.~\eqref{eq:dispatch}.
\end{tabular}
\end{center}
The kernel can consume the resulting certificate for $P=A$ without knowing how
$B$ was proposed. A transformation may supply a proof directly or submit an
untrusted candidate and witness to a verified checker.

Reusable condition services sit above the kernel. For example,
\code{sequence_readonly_conditions} composes a first check on $D$ with a second
check certified on $D\land P_1$. A refused first stage skips the second;
acceptance of both establishes $P_1\land P_2$. This stronger continuation domain
allows earlier checks to license later reads. Prefix-scan and probe-simplification
libraries expose further proof obligations rather than adding concrete memory
or optimizer semantics to the kernel.

\subsection{Local Reasoning and Complete Programs}

Language services establish local source/candidate execution bridges and the
frame that restores public observations. A domain implementation proves the
relation between the localized source and candidate models. Their composition
can discharge conditional equivalence, as implemented by
\code{localized_conditional_equivalence} in the read-only frontend.

Installation into a complete program is a separate language/IR obligation.
The generic \code{context_certificate} is a hook for a supplied lifting theorem:
it requires legal placement and admissibility of the actual replacement. It
does not prove that an arbitrary context preserves an arbitrary local relation.
The concrete Clight hosts establish control, private-resource, and progress
conditions and then compose with the CompCert backend.

Repeated rewrites are checked against their actual intermediate programs.
The read-only frontend's \code{certified_rewrite_sequence_equivalent} composes
finite sequences of installed equivalences. State relations, entry domains,
and placement evidence cannot simply be reused after an intervening rewrite
without showing that their requirements still hold.

\begin{table}[t]
\caption{Responsibilities at the proof boundary.}
\label{tab:responsibilities}
\centering\small
\begin{tabular}{>{\raggedright\arraybackslash}p{.19\textwidth}
  >{\raggedright\arraybackslash}p{.35\textwidth}
  >{\raggedright\arraybackslash}p{.36\textwidth}}
\toprule
Party&Supplies once or as a reusable library&Supplies for a transformation/site\\
\midrule
Kernel&Certificate interfaces and local composition&Consumes certificates for actual code/checks\\
Language/IR&Execution, primitive safety, dispatch, frames, installation laws&Site evidence satisfying host contracts\\
Domain implementer&Candidate checker, range/footprint and assumption services&Source/model bridge, sufficient premise, candidate and placement data\\
\bottomrule
\end{tabular}
\end{table}
